\documentclass[11pt,a4paper]{article}
\usepackage{fontspec}
\usepackage[a4paper,margin=2.2cm]{geometry}
\usepackage{booktabs,longtable,array,xcolor,enumitem,tcolorbox}
\usepackage{xurl}
\usepackage[hidelinks]{hyperref}
\setmainfont{Sylfaen}
\setmonofont{Consolas}
\newfontfamily\symfont{Segoe UI Symbol}
\setlength{\parindent}{0pt}\setlength{\parskip}{4pt}
\begin{document}
\section*{Colón-González et al. (2021): რა წაიკითხო და რა არის მნიშვნელოვანი}

\textbf{სტატია:} Colón-González, F. J., Soares Bastos, L., Hofmann, B., … Brady, O. J., Lowe, R. (2021).
\emph{Probabilistic seasonal dengue forecasting in Vietnam: A modelling study using superensembles.}
PLOS Medicine 18(3): e1003542. \url{https://doi.org/10.1371/journal.pmed.1003542} (CC BY 4.0)
\textbf{ფაილი:} \texttt{colon\_gonzalez\_2021\_plosmed.pdf} (30 გვერდი; PDF-ის გვერდი = ჟურნალის გვერდი, „N / 30“)

\begin{tcolorbox}[colback=blue!4,colframe=blue!40!black,boxrule=0.4pt]\textbf{საინტერესოა:} ბოლო ავტორი \textbf{Rachel Lowe}-ა, იგივე, ვინც ბრაზილიის სტატიის (Lowe et al. 2021) პირველი ავტორია.\end{tcolorbox}

\medskip\hrule\medskip

\section*{რა წაიკითხო (≈ 6 გვერდი, დანარჩენი გამოტოვე)}

\begin{longtable}{@{}p{0.317\linewidth}p{0.317\linewidth}p{0.317\linewidth}@{}}
\toprule
\textbf{გვერდი} & \textbf{სექცია} & \textbf{რატომ} \\
\midrule
6 & ბოლო აბზაცი Introduction-ში + \textbf{„Dengue surveillance data“} & შემთხვევების წყარო \\
7 & \textbf{„Historical demographic and land cover data“}, \textbf{„Historical Earth observation data“} & მოსახლეობა, კლიმატი, აგრეგაცია \\
8 & ზედა ნაწილი: \textbf{საზღვრების ცვლილება} & ძალიან მნიშვნელოვანი ჩვენთვის \\
9 & \textbf{„Model specification“} ბოლო აბზაცი & lag-ები და ინსტრუმენტები \\
1 & \textbf{Data Availability Statement} & ღიაა თუ არა მონაცემები \\
\bottomrule
\end{longtable}

მოდელის მათემატიკა (ფორმულები გვ. 8–11, superensemble, CRPS) \textbf{არ გჭირდება}: შენი თეზისი მონაცემების მომზადებაზეა.

\medskip\hrule\medskip

\section*{რას ამბობს ეს ნაწილი (ქართულად)}

\subsection*{1. შემთხვევები (გვ. 6)}
\begin{itemize}[nosep]
\item თვიური დენგეს შემთხვევები ვიეტნამის \textbf{ჯანდაცვის სამინისტროდან}, \textbf{2002-08 – 2019-12}, \textbf{63 პროვინცია}.
\item „Data comprised suspected and confirmed dengue cases“, მაგრამ არ წერია, რამდენია რომელი. სეროტიპები არ არის.
\item გვ. 7-ზე: 2002–2019 წლებში შემთხვევის განსაზღვრება \textbf{არ შეცვლილა}. $\rightarrow$ ეს \textbf{შედარებადობის} საკითხია (ჩვენი A20 / Clarke).
\end{itemize}

\subsection*{2. მოსახლეობა (გვ. 7)}
\begin{itemize}[nosep]
\item პროვინციის მოსახლეობა \textbf{SEDAC GPW v4.11}-დან, 1 კმ, \textbf{ყოველ 5 წელში} (2000–2020).
\item „Intervening years were generated using linear interpolation.“ $\rightarrow$ \textbf{ზუსტად ჩვენი FR-17} (წლიური $\rightarrow$ შუალედური წლები წრფივად).
\end{itemize}

\subsection*{3. კლიმატი (გვ. 7)}
ბევრი სხვადასხვა წყარო, სხვადასხვა რეზოლუციით და \textbf{წყაროს შეცვლა დროში}:
\begin{itemize}[nosep]
\item ტემპერატურა (min, max, mean): MODIS, 1 კმ;
\item ნალექი: TRMM 2014-04-მდე, \textbf{შემდეგ GPM} (წყარო შეიცვალა!);
\item ტენიანობა: MODIS;
\item ქარი: \textbf{ERA5} 2002–2011, \textbf{შემდეგ NOAA CFS} (ისევ წყაროს ცვლილება);
\item El Niño-ს ინდექსი: NOAA.
\end{itemize}
$\rightarrow$ \textbf{მოტივაციისთვის კარგი არგუმენტი:} ერთ კვლევაში 6+ წყარო, ზოგი ცვლადისთვის დროში ორი წყარო ერთმანეთზე „გადაბმული“. ეს ზუსტად ის პრობლემაა, რისთვისაც ჩვენი \textbf{წყაროების შედარება (FR-27)} არსებობს.

\subsection*{4. სივრცითი აგრეგაცია (გვ. 7–8): ყველაზე საინტერესო}
\begin{itemize}[nosep]
\item ბადის მონაცემები პროვინციებზე Python-ის მოდულით  გასაშუალოვდა; ნაწილობრივ დაფარულ უჯრაზე, წილის მიხედვით (როგორც Lowe-ს \texttt{exactextractr}).
\item \textbf{„Climate data were aggregated at the province level using population weighted averages for each month.“}
$\rightarrow$ კლიმატი \textbf{მოსახლეობით შეწონილი} საშუალოა: სადაც მეტი ადამიანი ცხოვრობს, ის უჯრა მეტს იწონის.
Lowe \textbf{ფართობით} წონიდა, ესენი \textbf{მოსახლეობით}. ეს მესამე ვარიანტია ჩვენი A/B/C-ს გვერდით.
\item წონებისთვის \textbf{WorldPop} (100 მ), რადგან SEDAC-ს რაიონის შიგნით ერთნაირი მნიშვნელობები ჰქონდა. WorldPop
მხოლოდ 2009, 2010, 2015, 2020 წლებისთვის იყო $\rightarrow$ შუალედური წლები \textbf{წრფივად ინტერპოლირდა}.
\end{itemize}

\subsection*{5. საზღვრების ცვლილება (გვ. 8): ჩვენთვის ყველაზე მნიშვნელოვანი}
\begin{itemize}[nosep]
\item „There were multiple changes in the administrative boundaries of the provinces over the study period. The last
change took place in 2008. Consequently, we froze the administrative boundaries of the country at their 2008 level.“
\item შემთხვევები 2008 წლის საზღვრებზე \textbf{რაიონების ჩანაწერებით} გადაანაწილა ჯანდაცვის სამინისტრომ.
\end{itemize}
$\rightarrow$ \textbf{იგივე პრობლემა, რაც ტაილანდის Bueng Kan-ში}, ოღონდ სხვა გადაწყვეტით:
\begin{itemize}[nosep]
\item Suttirat: ახალი პროვინცია ძველს \textbf{შეუერთეს};
\item Colón-González: საზღვრები \textbf{ერთ წელზე „გაყინეს“} (2008).
ორივე ერთსა და იმავე ნაბიჯს ასრულებს: \textbf{დროში ცვალებადი ერთეულების ჰარმონიზაცია}. ეს ახლა \textbf{სამ სტატიაშია}
(+ ვიეტნამის 2025 წლის გაერთიანება HDX-ზე) $\rightarrow$ FR-28-ის ძლიერი დასაბუთება.
\end{itemize}

\subsection*{6. lag-ები (გვ. 9)}
\begin{itemize}[nosep]
\item ტემპერატურა, ნალექი, ტენიანობა, DTR: \textbf{3-თვიანი მოძრავი საშუალო lag 0–2 თვე}.
\item El Niño: lag 0–3 თვე. ქარი: lag-ის გარეშე.
\item ასევე \textbf{წინა თვის შემთხვევები} (lag 1 შემთხვევებისთვის).
\item გვ. 4: „Most of the effects of climate on dengue incidence are delayed by 1 to 2 months“.
\end{itemize}

\subsection*{7. ინსტრუმენტები (გვ. 7, 9)}
\begin{itemize}[nosep]
\item Python 3.6 (\texttt{rasterstats}), R 3.6.1 + \textbf{INLA}. კოდი GitHub-ზე: github.com/FelipeJColon/paper\_dengue\_superensemble.
\end{itemize}

\subsection*{8. მონაცემების ხელმისაწვდომობა (გვ. 1)}
\begin{itemize}[nosep]
\item დენგეს მონაცემები: \textbf{მოთხოვნით}, ვიეტნამის ჯანდაცვის სამინისტროს კონკრეტული პირისგან. $\rightarrow$ \textbf{ღია არ არის}.
\end{itemize}

\medskip\hrule\medskip

\section*{მოტივაციისთვის გამოსადეგი წინადადებები (გვ. 4–5)}
\begin{itemize}[nosep]
\item გვ. 4: „Some of these challenges include the lack of multi-decadal health data sets …, **the common mismatch of
scales between climate data outputs and data used for decision-making** …“ $\rightarrow$ მასშტაბების შეუსაბამობა = ჩვენი
ჰარმონიზაცია.
\item გვ. 5: „… difficulties related to the integration of these systems into existing public health procedures and
\textbf{occasional lack of technical expertise}.“ $\rightarrow$ უნარების ნაკლებობის კიდევ ერთი (ფრთხილი) წყარო.
\end{itemize}

\medskip\hrule\medskip

\section*{ჩვენი ცხრილისთვის (Schritte / Traceability)}

\begin{longtable}{@{}p{0.237\linewidth}p{0.237\linewidth}p{0.237\linewidth}p{0.237\linewidth}@{}}
\toprule
\textbf{ნაბიჯი} & \textbf{Colón-González 2021} & \textbf{გვ.} & \textbf{ჩვენი სისტემა} \\
\midrule
თვიური შემთხვევები & MoH, 63 პროვინცია, 2002–2019 & 6 & {\symfont ✓} (FR-28 ნაწილი 1, ტაილანდი) \\
საზღვრების ჰარმონიზაცია & საზღვრები 2008-ზე გაყინული & 8 & {\symfont ✓} ტაილანდი (Bueng Kan), {\symfont {\symfont ⚠}} ზოგადად \\
კლიმატის რამდენიმე წყარო, წყაროს ცვლილება დროში & TRMM$\rightarrow$GPM, ERA5$\rightarrow$CFS & 7 & {\symfont {\symfont ⚠}} შედარება {\symfont ✓} (FR-27), გადაბმა {\symfont ✗} \\
რასტრი $\rightarrow$ პროვინცია & rasterstats, მოსახლეობით შეწონილი & 7 & {\symfont ✗} (FR-28 ნაწილი 2, ვარიანტები A/B) \\
მოსახლეობა, შუალედური წლები & GPW 5-წლიანი, წრფივი ინტერპოლაცია & 7 & {\symfont ✓} (FR-17) \\
lag-ები & 3-თვიანი საშუალო lag 0–2; ENSO 0–3 & 9 & {\symfont ✓} lag სვეტები (FR-23); მოძრავი საშუალო {\symfont ✗} \\
ინსტრუმენტები & Python + R/INLA & 7, 9 & — \\
\bottomrule
\end{longtable}

\medskip\hrule\medskip

\section*{შენი დავალება}
\begin{enumerate}[nosep]
\item გახსენი PDF და \textbf{ზემოთ მოყვანილი ციტატები გადაამოწმე} (Ctrl+F), გვერდები ჩამიწერე, თუ რამე არ ემთხვევა.
\item შეავსე \texttt{kitxvis\_furceli.md}-ის ასლი (\texttt{kitxvis\_furceli\_colon\_gonzalez\_2021.md}) \textbf{შენი სიტყვებით} — ეს
შეჯამება დახმარებაა და არა შენი ფორმის ჩანაცვლება.
\item ნაწილი 7, „ჩემი დასკვნა“: რას ამატებს ეს სტატია? (მინიშნება: საზღვრები, მოსახლეობით შეწონვა, წყაროს ცვლილება დროში.)
\end{enumerate}
\end{document}
